%! Systems Involving Quadratics — Unit 3, Lesson 3.6 — Exit Ticket (Answer Key)
\documentclass[10pt]{article}
\usepackage{algebra2-article}
\usepackage{algebra2-key}   % pulls in -boxes; do NOT also load -boxes
\newcommand{\TallMath}[1]{$\displaystyle #1\rule[-1.4em]{0pt}{3.2em}$}

\begin{document}

\pageheader{Unit 3, Lesson 3.6}{Exit Ticket --- Answer Key}
\namedateperiod

\vspace{0.12in}

No notes. Show your steps. A solution is an ordered pair --- find \emph{both} coordinates and check them.

\begin{enumerate}[leftmargin=1.8em, itemsep=14pt]
  \item \textbf{Solve the system by substitution.} Write each solution as an ordered pair.
        \[ y=x^2-x-1, \qquad y=x+2 \qquad\Rightarrow\qquad \ans{$(3,5)$ and $(-1,1)$} \]
        \ansline{Set equal: $x^2-x-1=x+2\Rightarrow x^2-2x-3=0\Rightarrow(x-3)(x+1)=0\Rightarrow x=3,-1$. Back-sub into $y=x+2$: $x=3\to5$, $x=-1\to1$.}

  \item \textbf{Number of solutions --- no solving.} How many solutions does the system
        $\ y=x^2+2,\ \ y=2x\ $ have? Substitute, then use the discriminant to justify.
        \[ b^2-4ac=\ans{$-4$} \qquad\Rightarrow\qquad \ans{$0$ solutions (the line misses the parabola)} \]

  \item \textbf{SOL-style multiple choice.} After substitution, a linear--quadratic system collapses to
        $x^2-6x+9=0$. How many solutions does the system have?
        \begin{enumerate}[label=\Alph*., leftmargin=1.9em, itemsep=1pt]
          \item $0$ solutions
          \item exactly $1$ solution \textcolor{keyred}{\textbf{$\leftarrow$ correct}}
          \item exactly $2$ solutions
          \item infinitely many solutions
        \end{enumerate}
        \begin{itemize}[leftmargin=1.4em, nosep]
          \item[] $x^2-6x+9=(x-3)^2=0$ has the double root $x=3$ ($b^2-4ac=0$), so the line is \emph{tangent} to the parabola --- exactly one shared point.
        \end{itemize}
\end{enumerate}

\begin{teachernote}
Sort into ``got it / found $x$ but no $y$ / mismatched coordinates / didn't check both.'' Telltales: item~1 stopping
at $x=3,-1$ without producing pairs, or back-substituting into the parabola and slipping; item~2 a sign error making
$b^2-4ac\ge0$, or reading ``$-4$'' as ``no solution of any kind'' (here it genuinely \emph{is} $0$ real solutions,
since a system's solutions must be real intersection points); item~3 choosing \textbf{C} (the ``a parabola meets a
line twice'' default) instead of recognizing the $D=0$ tangent case. If item~3 or the back-substitution step is
widely missed, reopen Lesson~3.7 with a 60-second ``a solution is a \emph{pair}, and $D$ counts intersections'' recap.
\end{teachernote}

\end{document}
